% =====================================================================
% VERSION INSTITUT A 150 PAGES : annexe des adaptations par pilier et du cas d'usage.
% Reprise de supplement/S1_adaptations_pilier.tex (ancienne annexe C de la version Institut de
% 195 pages), reecrite en phrases courtes. Les nombres, formules, etiquettes et references sont
% ceux de l'original. Lue par main_institut_150.tex seul.
% =====================================================================

\chapter{Adaptations par pilier et cas d'usage}
\label{chap:adaptations}

La queue de sévérité, l'accumulation par les tiers et la détection appellent des adaptations
distinctes. Les diagnostics suivants en comparent les effets sous des hypothèses déclarées ; ils
ne calibrent pas les mécanismes par pilier. Le classement peut changer selon l'assignation des queues.

\section{Queue de sévérité par pilier}
\label{sec:xi-pilier}

Une seule queue $\xi$ est appliquée aux cinq piliers. La queue est pourtant hétérogène selon la
catégorie Bâle. Sur la table recalculée pour ce travail, $\xi$ varie de $0{,}87$
(\emph{Employment}) à $1{,}37$ (\emph{Execution, Delivery \& Process Management}), soit une
dispersion de $0{,}50$.
% SOURCES-SCRIPTS: 67
L'hétérogénéité est un fait, son assignation aux piliers n'en est pas un. Les catégories Bâle
décrivent la nature de la perte, pas les domaines de contrôle de DORA, et la plus proche des TIC
(\emph{Business Disruption \& System Failures}) ne compte que $111$ observations, sous le seuil
d'estimation de $150$. Nous ne posons donc pas $\xi_j = f(\text{pilier})$. Nous suivons la logique du
chapitre~\ref{chap:partielle} : chaque pilier reçoit l'une des cinq queues observées, nous énumérons
les $120$ assignations, et le besoin de capital est borné. Pour isoler l'effet de
l'hétérogénéité, les cinq queues sont centrées sur la valeur de référence.

Même à moyenne des queues fixe, donner à chaque pilier sa propre queue relève le besoin de capital de $+55\,\%$ à
$+96\,\%$. Cet effet tient à la convexité. La VaR croît plus vite que linéairement en $\xi$, de sorte qu'une
queue lourde sur un pilier n'est pas compensée par une queue légère ailleurs. Une queue commune
sous-estime ainsi le capital, ce qui justifie de lire le niveau en bande. Le classement
résiste : le pilier de tête reste dominant dans $83\,\%$ des $30$ assignations tirées parmi les $120$, sans être invariant
(figure~\ref{fig:xi-pilier}).

\begin{figure}[H]\centering
\figover{Z8_xi_par_pilier.png}
\caption{Queue de sévérité par pilier. (a)~Besoin de capital borné sur les $120$ assignations des queues
observées aux piliers ; une queue commune (trait) est en bas de la bande. (b)~Le pilier de tête peut changer selon l'assignation ; les positions affichées sont celles du diagnostic.}
% SOURCES-SCRIPTS: 37
\label{fig:xi-pilier}
\end{figure}

\section{Pilier des tiers et accumulation}
\label{sec:p4-accumulation}

P4 n'est pas un n\oe{}ud de cascade comme les autres. La défaillance d'un prestataire partagé dégrade
plusieurs piliers à la fois, par dépendance commune, au lieu de se propager de proche en proche.
C'est le mécanisme de MOVEit, où $191$ victimes tombent d'un seul choc \matcomp.
% SOURCES-SCRIPTS: 08g
Nous lui substituons un \emph{choc commun} : lorsque l'incident commence en P4, chaque autre pilier
est inclus indépendamment avec une probabilité $\omega$, au lieu d'une propagation dirigée.

Le choc commun produit des co-occurrences sans direction. Les incidents chez un prestataire partagé
motivent ce mécanisme, mais ne calibrent ni son intensité au sein d'une entité ni la partie symétrique
$S$ de la décomposition $W=S+A$. Celle-ci reste posée à dire d'expert
(chapitre~\ref{chap:partielle}). La concentration du nuage sert d'ancrage au scénario ; elle ne mesure
pas une probabilité de transmission entre piliers.

L'intensité $\omega$ n'est pas identifiée au sein d'une entité isolée, car le \emph{batch}
observé est un phénomène de portefeuille. Nous la bornons sur $[0,\gamma]$, où $\gamma=0{,}68$ est la
concentration du nuage déjà retenue. Le besoin de capital se lit alors en bande. À $\omega=\gamma$, une
dépendance partagée emporte en moyenne $3{,}7$ piliers d'un coup, soit $+19\,\%$ de besoin de capital par rapport
au traitement en n\oe{}ud de cascade.

Dans ce scénario, le choc commun augmente la fréquence des
incidents multi-piliers : la probabilité qu'un incident P4 touche au moins trois piliers passe de
$0{,}22$ avec la cascade à $0{,}90$ avec le choc commun (figure~\ref{fig:p4}). Ce calcul compare les
ensembles atteints au cours d'un incident, sans estimer les délais entre défaillances.

\begin{figure}[H]\centering
\figover{Z9_p4_accumulation.png}
\caption{Le pilier des tiers en accumulation. (a)~Besoin de capital borné sur l'intensité du choc partagé
$\omega\in[0,\gamma]$, comparé au traitement de P4 en n\oe{}ud de cascade (trait).
(b)~Probabilité de toucher plusieurs piliers au cours d'un incident amorcé en P4, selon le
mécanisme retenu. Le caractère commun du choc ne suffit pas à identifier son intensité.}
% SOURCES-SCRIPTS: 38
\label{fig:p4}
\end{figure}

\subsection{Seuil de déclenchement et dépendance de queue}
\label{sec:p4-seuil}

Deux questions se séparent ici : celle de la marge, et celle de la dépendance. Le déclenchement d'un
pilier est un franchissement de seuil,
\begin{equation}
  \{\text{pilier } j \text{ déclenché}\} \;=\; \{X_j \ge K_j\},
  \qquad K_j = F^{-1}(1-p_j),
  \label{eq:seuil-Kj}
\end{equation}
et pour $P_4$, $p_4=0{,}151$ donne $K_4=1{,}03$. La loi $F$ de la latente n'y sert que de
normalisation. Par construction, $\Prob(X_j\ge K_j)=p_j$ quelle que soit $F$ continue et strictement
croissante, donc le choix de $F$ ne porte aucune information sur la fréquence de déclenchement. La
queue lourde du capital se trouve dans la sévérité (GPD, chapitre~\ref{chap:socle}), pas dans ce
seuil. Une marge normale est une normalisation à copule et probabilités fixées, pas une validation du modèle.

L'accumulation est gouvernée par la dépendance de queue supérieure de la copule $C$ qui lie les
latentes,
\begin{equation}
  \lambda_U \;=\; \lim_{u\to 1^-}\Prob\bigl(U_2>u \,\big|\, U_1>u\bigr)
             \;=\; \lim_{u\to 1^-}\frac{1-2u+C(u,u)}{1-u} .
  \label{eq:lambda-U}
\end{equation}
Une copule gaussienne non dégénérée vérifie $\lambda_U=0$ pour une corrélation strictement inférieure à un. La copule de Gumbel de
paramètre $\theta>1$ donne
\begin{equation}
  \lambda_U^{\text{Gumbel}} \;=\; 2-2^{1/\theta},
  \label{eq:lambda-gumbel}
\end{equation}
soit $0{,}53$ pour le $\theta=1{,}8$ retenu ailleurs dans ce travail. Retenir une copule gaussienne
ici sous-estimerait donc le co-extrême, en contradiction avec ce choix.

Nous ne tranchons pas la copule. Nous lisons le co-déclenchement en bande, à $\tau$ de Kendall égal,
entre le plancher gaussien et le haut de Gumbel. La fréquence marginale d'un pilier reste
$p_4=0{,}15$ dans les deux cas, car la copule ne change pas combien un pilier bascule seul, mais s'ils basculent
\emph{ensemble}. Seul le co-extrême diverge : la probabilité que les cinq piliers tombent avec P4 passe
de $0{,}17$ pour la gaussienne à $0{,}39$ pour Gumbel, soit un facteur $2{,}3$. C'est le scénario
MOVEit que la queue gaussienne manque (figure~\ref{fig:p4-seuil}). Le co-déclenchement d'un autre
pilier tient dans $[0{,}48\,;\,0{,}58]$, sous la borne $\gamma=0{,}68$ du choc commun, qui reçoit
ainsi une origine structurelle. Tout est exprimé par pilier, c'est-à-dire en nombre de piliers de l'entité plutôt qu'en fraction de
portefeuille. La vue inter-clients, avec les centaines d'entités victimes de MOVEit, relèverait d'une
loi de Vasicek en dimension entité et sort du besoin de capital par entité.

\begin{figure}[H]\centering
\figover{Z12_seuil_declenchement_p4.png}
\caption{Le seuil de P4 et l'accumulation bornée. (a)~La marge : franchissement de $K_4$, où la
normale n'est qu'une normalisation. (b)~À dépendance moyenne égale, le co-dépassement s'effondre vers
$0$ pour la gaussienne mais plafonne à $0{,}53$ pour Gumbel. (c)~L'accumulation par pilier, Gumbel
concentrant la masse sur les cinq piliers groupés.}
% SOURCES-SCRIPTS: 42
\label{fig:p4-seuil}
\end{figure}

\subsection{Déclenchement comme temps d'arrêt}
\label{sec:p4-dynamique}

Le seuil $z^*$ est jusqu'ici un niveau statique. Si le facteur tiers est un processus continu $(Z_t)_{t\ge0}$ adapté à une filtration $(\mathcal{F}_t)$, le déclenchement devient un temps de premier passage,
\begin{equation}
  \tau \;=\; \inf\{\,t\ge 0 \;:\; Z_t \ge z^*\,\},
  \label{eq:temps-arret}
\end{equation}
qui est un temps d'arrêt, puisque $\{\tau\le t\}=\{\sup_{s\le t}Z_s\ge z^*\}\in\mathcal{F}_t$.

\begin{proposition}[Loi du délai de déclenchement]\label{prop:premier-passage}
Si $(Z_t)$ est un mouvement brownien standard issu de $0$ et $z^*>0$, alors pour tout horizon $T>0$
\begin{equation}
  \Prob(\tau\le T) \;=\; 2\,\Phi\!\left(-\frac{z^*}{\sqrt{T}}\right)
  \;=\; 2\,\Prob(Z_T\ge z^*).
  \label{eq:reflexion}
\end{equation}
\end{proposition}

\begin{proof}
C'est le principe de réflexion. Par symétrie du mouvement brownien après $\tau$, on a
$\Prob(\tau\le T,\;Z_T< z^*)=\Prob(\tau\le T,\;Z_T> z^*)=\Prob(Z_T>z^*)$, la dernière égalité tenant
car $\{Z_T>z^*\}\subset\{\tau\le T\}$. En sommant les deux événements disjoints,
$\Prob(\tau\le T)=2\,\Prob(Z_T>z^*)=2\,\Phi(-z^*/\sqrt{T})$.
\end{proof}

La lecture dynamique est le double de la lecture statique $\Prob(Z_T\ge z^*)$, puisqu'elle compte tous
les franchissements survenus avant l'horizon, et pas seulement la position finale. À $\rho=0{,}50$, la
probabilité de déclenchement de P4 en moins d'un an est de $0{,}14$, pour un délai médian de l'ordre
de cinq ans. À $\rho=0{,}75$, elle monte à $0{,}23$, pour un médian d'environ trois ans
(figure~\ref{fig:p4-dynamique}).

\begin{figure}[H]\centering
\figover{Z16_p4_temps_arret.png}
\caption{Le déclenchement de P4 comme temps d'arrêt. (a)~Le facteur tiers (martingale) et le premier
passage de $z^*$. (b)~La loi du délai de déclenchement. (c)~La lecture dynamique
$\Prob(\tau\le 1\text{ an})$ vaut environ le double de la statique.}
% SOURCES-SCRIPTS: 52 128
\label{fig:p4-dynamique}
\end{figure}

Le principe de réflexion utilise la propriété de Markov forte et la symétrie des accroissements
du brownien ; la propriété de martingale seule ne suffit pas. La forme de Dirichlet du
chapitre~\ref{ch:identifiabilite} est un objet distinct. Aucun changement de mesure n'intervient.
Le cadre complet est déporté \matcomp.
En revanche, le comptage journalier des brèches du secteur financier (2016--2023) se diagnostique par son
résidu compensé, qui est une différence de martingale si l'intensité prévisible est bien spécifiée.

Une
tendance seule laisse une autocorrélation de $+0{,}26$ au premier retard ; une excitation des jours
antérieurs la ramène à $-0{,}02$, avec un branchement de $0{,}29$, puis de $0{,}09$ sans la semaine MOVEit.
Les seuils sont recalibrés par bootstrap avec réestimation et la spécification reste rejetée. Les
vulnérabilités du catalogue CISA KEV n'améliorent pas l'ajustement sur 2022--2023 ($p=0{,}90$), et la
faille MOVEit y est ajoutée quatre jours après son pic. Ce diagnostic ne prouve ni causalité ni contagion.
Le brownien pur atteint le niveau presque sûrement, mais son délai est d'espérance infinie. Un facteur
à retour à la moyenne (Ornstein-Uhlenbeck) donnerait un taux de déclenchement récurrent et fini, plus
réaliste pour un facteur systémique. Nous gardons le brownien parce qu'il permet un calcul explicite.

\subsection{Sous-processus tiers sur donnée ouverte}
\label{sec:p4-sousprocess}

Descendre P4 au niveau du sous-processus (quel sous-traitant, quelle fonction critique, quelle
substituabilité) suppose le registre d'externalisation de DORA. Sa structure est fixée par la norme
technique d'exécution \citep{eu2024_2956}, mais seules les données de l'entité le renseignent. La
donnée ouverte tranche déjà deux points opposés (figure~\ref{fig:p4-open}).

La taxonomie fine du tiers y est trop rare. La sous-catégorie \og Vendors \& Suppliers \fg{} d'OpRisk
ne compte que $20$ incidents en finance, soit $1\,\%$ des incidents TIC, et \texttt{actor.Partner} du
VCDB en compte $56$. Ces effectifs tombent à zéro deux niveaux plus bas, ce qui renvoie au registre.
La concentration et l'origine tierces sont en revanche ancrées sur des données européennes
officielles, qui permettent de situer deux paramètres jusque-là posés.
\begin{itemize}\itemsep2pt
\item Le premier rapport d'incidents DORA des autorités européennes \citep{esas2026_incidents}
      établit que $29\,\%$ des $3\,383$ incidents TIC majeurs de $2025$ proviennent d'une défaillance
      de tiers. P4 est donc une amorce majeure, ce que le classeur posait par jugement
      ($\text{\texttt{ROOT}}[4]=0{,}90$).
\item L'analyse des registres d'externalisation par la BCE \citep{ecb2024_outsourcing} montre que plus
      de $30\,\%$ du budget des grandes banques va à $10$ fournisseurs, et $50\,\%$ du budget critique
      à $30$. Les autorités européennes ont désigné $19$ fournisseurs TIC critiques. Ces chiffres motivent le scénario d'accumulation, sans identifier $\omega$ ni mesurer la
charge de facteur $\gamma=0{,}68$. Concentration du marché et probabilité de choc sont distinctes.
\end{itemize}

\begin{figure}[H]\centering
\figover{Z14_p4_sousprocess_open.png}
\caption{Le sous-processus tiers sur donnée ouverte. (a)~La taxonomie fine du tiers est sous le
seuil d'estimation (registre DORA requis). (b)~La concentration est nette, côté exposition (BCE)
comme côté sinistralité (Gini $0{,}69$). (c)~Origine et concentration tierces ancrées sur donnée
européenne officielle.}
% SOURCES-SCRIPTS: 44 08g
\label{fig:p4-open}
\end{figure}

\section{Détection de P2 et P3 comme courbe de survie}
\label{sec:detection}

P2 (gestion des incidents) et P3 (tests) ne sont pas des n\oe{}uds de sévérité. Ils gouvernent le
temps de détection, car un incident non détecté à temps escalade et devient matériel. Le modèle
résume cet effet par un multiplicateur du taux de matérialité par pilier touché, de $0{,}85$ à l'état conforme et
$1{,}20$ à l'état non conforme. Nous le remplaçons par une lecture de survie ancrée sur donnée.

Le VCDB donne la distribution du délai de découverte en finance ($n=278$). Elle est bimodale :
$47\,\%$ des incidents sont découverts le jour même, $8\,\%$ en semaines et $45\,\%$ en mois ou plus.
La matérialité vient des incidents à découverte lente, d'où un taux de matérialité par pilier touché proportionnel à
la part de découverte lente, ancrée à $\pi_{\text{slow}}=0{,}45$.

Ce point est mesuré, mais l'élasticité de la détection à la maturité de P2 et P3 (de combien un
pilier conforme détecte-t-il plus vite) ne l'est pas, et nous la bornons. Le surcoût du seul canal
détection est alors une bande, de $1\,775$ à $8\,052$~M\euro{} selon l'élasticité. Le multiplicateur
grossier $0{,}85/1{,}20$ est retrouvé pour une élasticité $\eta\approx 0{,}36$. Il est donc compatible avec un point de la
courbe de scénario ancrée sur le VCDB, sans que son élasticité soit estimée (figure~\ref{fig:detection}).

\begin{figure}[H]\centering
\figover{Z10_detection_survie.png}
\caption{P2 et P3, la détection comme courbe de survie. (a)~Distribution bimodale observée au VCDB.
(b)~Le multiplicateur de matérialité comme point de la courbe de survie ; le multiplicateur grossier
actuel correspond à $\eta\approx0{,}36$. (c)~Surcoût du canal détection, borné sur l'élasticité.}
% SOURCES-SCRIPTS: 39 67
\label{fig:detection}
\end{figure}

\section{Cas d'usage, des adaptations aux indicateurs pilotables}
\label{sec:kpi}
% SOURCES-SCRIPTS: 43 68 82 88 74

Les adaptations se lisent comme un tableau de bord opposable à un souscripteur ou à un superviseur.
Chaque indicateur observable est rattaché à un canal du moteur : fréquence d'entrée, délai de
notification (P2), délai de confinement (P3), concentration des tiers (P4). Le capital en jeu derrière
chacun se mesure par l'écart entre l'état conforme et l'état non conforme, soit $14\,139$~M\euro{}
($6\,049 \to 20\,188$) pour le périmètre financier étudié, avec un écart-type de $734$ entre quatre graines.

\subsection{Trois lectures d'un canal}
% SOURCES-SCRIPTS: 43 68 82

On peut partir de l'état conforme et ne dégrader que ce canal (lecture \emph{isolée}), partir de
l'état non conforme et ne remédier qu'à ce canal (marginal de \emph{fermeture}), ou moyenner sur les
vingt-quatre ordres d'entrée (valeur de \emph{Shapley}). La figure~\ref{fig:interaction-canaux}(b)
compare les trois lectures, dont une seule somme à l'écart total. La lecture isolée manque
$5\,001$~M\euro{} et celle de fermeture dépasse l'écart d'autant, le rapport des deux allant de $1{,}5$
à $3{,}1$ selon le canal. Remédier à un canal rapporte donc davantage à une entité défaillante partout
qu'à une entité déjà conforme, puisque ce canal ferme aussi les effets croisés qu'il portait. La lecture de fermeture répond au scénario d'une entité défaillante sur tous les canaux.

Deux des trois sommes ne mesurent rien, et il ne faut pas les resommer. La question a été posée en
séance sous la forme d'un soupçon de double comptage. Il n'y a pas de double comptage, et c'est une
identité. En notant $m(S)$ les coefficients de Möbius,
\[
  \sum_i \text{isolé}(i) = \!\!\sum_{|S|=1}\!\! m(S), \qquad
  \sum_i \text{fermeture}(i) = \sum_S |S|\,m(S), \qquad
  \sum_i \varphi_i = \sum_S m(S) = \Delta .
\]
Quand on ferme un canal, il emporte tous les croisés auxquels il participe, et un croisé d'ordre $k$
est donc compté $k$ fois :
$1\times 9\,138 + 2\times 5\,345 + 3\times(-691) + 4\times 346 = 19\,141$, à la précision machine, aux arrondis d'affichage près.
L'encadrement $9\,138 \leq 14\,139 \leq 19\,141$ est en revanche constaté ici, non garanti : la
sur-modularité de la fonction de capital l'assurerait, mais la paire propagation\,$\times$\,accumulation ($-330$) y échappe.

Shapley est la seule lecture qui
répartisse l'interaction sans la perdre ni la dupliquer. Elle reste une convention d'attribution, pas
une prévision, car deux des quatre canaux sont bornés et leur part n'est pas un budget actionnable. Le défaut
de la colonne isolée est celui que rencontre l'analyse de sensibilité lorsque les entrées cessent
d'être indépendantes : les indices de Sobol ne somment plus au total, alors que la valeur de Shapley y
somme par construction \citep{Owen2014, IoossPrieur2019}. Ces travaux décomposent toutefois une
variance, alors que la fonction caractéristique employée ici est un écart de capital. L'argument se
transpose, pas le résultat.

\subsection{L'état non conforme comme coin supérieur}
\label{sec:coin-superieur}
% SOURCES-SCRIPTS: 88

Sur les soixante-cinq paires emboîtées des seize configurations, relâcher un canal de plus ne fait
jamais baisser le capital estimé, et cela vaut graine par graine, pas seulement en moyenne. Parmi ces
seize configurations, le quantile estimé le plus élevé est donc celui du coin haut, sur chacune des
graines utilisées, ce qui suffit à un test de résistance limité aux deux valeurs de chaque canal. Les valeurs intermédiaires des canaux ne sont pas testées, et
aucune démonstration n'étend ce constat au pavé continu. C'est la forme que prend ici le résultat de statique comparative monotone du préprint cité au
chapitre~\ref{ch:positionnement} \citep{ccrn2026}.

Deux réserves s'y attachent. D'abord, l'ordre observé porte sur les quantiles estimés, pas trajectoire par
trajectoire. À uniformes communs, la perte d'une année donnée baisse dans $16{,}2\,\%$ des cas lorsque l'on relâche la
propagation et dans $30{,}7\,\%$ lorsque l'on relâche la détection. Les motifs diffèrent : la table
des sous-ensembles dans le premier cas, la transformation de sévérité dans le second. La version
trajectorielle demanderait un couplage monotone, donc d'autres tirages, ce que le gel de la calibration
interdit pour renforcer un énoncé sans déplacer une conclusion. Ensuite, un coin peut être infaisable.
Deux des quatre canaux sont des bornes posées, et rien ne garantit qu'une entité présente les quatre
canaux à leur maximum en même temps. Le coin haut est le maximum observé parmi les seize
configurations évaluées, non la description d'une entité.

\subsection{Test de résistance inversé}
\label{sec:inverse}
% SOURCES-SCRIPTS: 68 84 88 92

L'ordre observé sur les configurations évaluées autorise la question réciproque, que pose aussi un dispositif ORSA. Au lieu de fixer un
état et de lire un capital, nous fixons le capital et cherchons les états qui le produisent.
Parmi ces configurations, celles qui atteignent une cible forment un ensemble croissant, décrit par
ses éléments minimaux ; pour les cibles étudiées, ils ne sont jamais plus de quatre sur seize.

\begin{center}\footnotesize
\renewcommand{\arraystretch}{1.2}
\begin{tabular}{@{}r l l@{}}
\toprule
\textbf{Cible} & \textbf{configurations minimales classées qui l'atteignent} & \textbf{minimum de canaux} \\
\midrule
$8\,000$ & fréq. ; dét.$+$prop. ; dét.$+$accum. ; prop.$+$accum. & $1$ \\
$10\,000$ & fréq. ; dét.$+$prop.$+$accum. & $1$ \\
$12\,000$ & fréq.$+$dét. ; fréq.$+$prop. ; fréq.$+$accum. & $2$ \\
$15\,000$ & les trois triplets contenant la fréquence & $3$ \\
$18\,000$ & les quatre canaux & $4$ \\
\bottomrule
\end{tabular}
\end{center}

\vspace{4pt}
\noindent Deux nombres résument le tableau. Le canal de fréquence seul atteint $10\,377$~M\euro, soit
$1{,}72$ fois l'état conforme, alors que les trois autres canaux réunis n'atteignent que $11\,413$,
soit $1{,}89$ fois. L'ordre de ces deux bornes tient sur les quatre graines. Toute cible sous le
premier seuil se produit donc par un canal unique, et toute cible au-dessus du second exige la
fréquence. Les trois autres canaux, pris seuls, restent sous la cible la plus basse : la détection plafonne à
$7\,497$~M\euro, la propagation à $7\,682$ et l'accumulation à $7\,777$. Dans ce modèle, être exposé à un
seul canal revient à être exposé à la fréquence.

Le canal d'accumulation n'admet pas d'inversion continue, et son effet isolé est un solde net. À
l'état conforme, le pilier des tiers est un n\oe{}ud de cascade ordinaire ; à l'état non conforme, il devient
un choc commun. Ce sont deux structures de table différentes, pas les extrémités d'un continuum. Un test inverse ne
peut donc pas demander à ce canal jusqu'où il doit se dégrader. La bascule à paramètre nul
fait baisser le capital de $239$~M\euro, résolu sur les quatre graines, parce qu'elle retire la
propagation propre du pilier avant d'ajouter le choc commun. L'effet isolé publié, $1\,728$~M\euro{},
est le solde d'un retrait de $239$ et d'un ajout de $1\,967$. On ne peut donc pas sommer les quatre
canaux.

Deux réserves restent à faire. Une cible déclarée inatteignable l'est parmi les seize configurations
évaluées, jamais pour une entité, puisque deux canaux sur quatre sont des bornes posées. Et une configuration proche d'une cible n'est pas
classée, parce que son appartenance dépend de la graine. Une frontière non résolue est d'ailleurs un
résultat : elle indique que la cible tombe dans le bruit de simulation.

\subsection{Décomposition de l'interaction}
% SOURCES-SCRIPTS: 68 84 88 100

L'interaction n'est plus un résidu, elle se décompose et vit aux paires. La décomposition de Möbius
sur le treillis des seize configurations rend compte des $+5\,001$~M\euro{} terme par terme. Elle compte quatre
effets principaux ($9\,138$~M\euro), six croisés de paires ($5\,345$), quatre croisés de triplets
($-691$) et un croisé des quatre canaux ($+346$). Leur somme vaut $14\,139$ à la précision machine, aux arrondis d'affichage près.
Cette identité atteste que la table est complète ; elle ne vérifie rien statistiquement.
Les trois paires porteuses passent toutes par la fréquence et font $86\,\%$ de l'ordre~2 :
fréquence\,$\times$\,détection $1\,739\pm708$, fréquence\,$\times$\,accumulation $1\,666\pm354$ et
fréquence\,$\times$\,propagation $1\,182\pm444$~M\euro. Ces $\pm$ sont des étendues de simulation à
paramètres fixés. Une loi d'échelle le montre : l'écart-type tombe de $1\,785$ à $339$~M\euro{}
lorsque le nombre d'années simulées passe de $10\,000$ à $160\,000$, soit une pente de $-0{,}60$ contre
$-0{,}50$ attendu pour du Monte-Carlo. Cette largeur se réduit en simulant davantage. L'incertitude
de paramètre est dix fois plus grande : le même croisé va de $429$ à $7\,998$~M\euro{} aux deux bornes
de l'IC90 de $\hat\xi$. Les deux premières paires ne sont séparées que de $73$~M\euro{}, et leur ordre s'inverse en perte
moyenne. On ne peut donc pas les classer.

Sur les seize configurations, quatorze mélangent des canaux conformes et non conformes. Ce sont
des contrefactuels de calcul pour attribuer l'écart. Une entité peut avoir des contrôles de maturités
différentes ; la donnée ne montre pas qu'elle puisse atteindre exactement les valeurs posées de
chaque configuration. Un état mixte doit donc rester présenté comme un scénario conditionnel.

Une seule paire est négative : propagation\,$\times$\,accumulation, à $-330\pm242$~M\euro, sur les
seize graines et les deux métriques. Les deux canaux se substituent l'un à l'autre au lieu de se
compléter. La cascade et le choc commun sont deux manières de faire tomber plusieurs piliers sur un même
sinistre, et quand l'une joue, il reste moins à gagner à l'autre. Le capital est donc super-additif en
euros, mais sous-additif entre ses deux canaux de co-occurrence. Cela borne l'amplitude que l'on peut
imputer à la contagion en général.

Les canaux se composent de façon quasi multiplicative. Une décomposition additive d'effets qui se
multiplient produit une interaction positive par simple arithmétique, $(1+a)(1+b)-1$ dépassant $a+b$.
La décomposition refaite sur le logarithme du capital le montre. Le produit des quatre facteurs isolés
vaut $3{,}47$ quand le facteur publié vaut $3{,}34$, soit $0{,}961$ fois ce produit. Un modèle
purement multiplicatif produirait $5\,813$~M\euro{} d'interaction, dont les $5\,001$~M\euro{} publiés
font $86\,\%$. Aucune des seize configurations ne s'écarte de plus de $8{,}1\,\%$ du produit de ses
facteurs isolés. L'interaction en euros s'explique donc par la multiplicativité, et le résidu propre est
une légère sous-multiplicativité, portée par la paire propagation\,$\times$\,accumulation.

Deux canaux suivent même une forme fermée. Par l'approximation de la perte unique \citep{bockerkluppelberg2005}, le
quantile croît comme $(\lambda\,p_u)^{\xi}$, ce qui prédit un facteur
$(53{,}62/21{,}56)^{0{,}5954}=1{,}720$ pour la fréquence contre $1{,}715$ simulé, et $1{,}228$ pour la
détection contre $1{,}239$.

Les conclusions de gestion restent les mêmes : remédier à un canal rapporte toujours davantage à
une entité défaillante partout. Leur statut change en revanche, car elles reposent sur une forme
fermée contrôlée par la simulation et non sur la seule simulation. L'interaction ne se lit donc pas
comme un effet propre de la contagion (figure~\ref{fig:interaction-canaux}).
% SOURCES-SCRIPTS: 100

\begin{figure}[H]\centering
\figover{S24_interaction_canaux.png}
\caption{La table complète des quatre canaux. (a)~Réconciliation par ordre : les quinze termes de
Möbius somment à l'écart, exactement. (b)~Les trois lectures d'un même canal, isolée, Shapley et
fermeture. (c)~Les six croisés de paires avec leur écart-type de simulation ; une seule est négative.}
\label{fig:interaction-canaux}
\end{figure}

\subsection{Classement des indicateurs par capital et par statut}
% SOURCES-SCRIPTS: 43 68

La fréquence et la détection pèsent lourd en capital et sont calibrables : ce sont les leviers de
remédiation à activer en premier. La propagation et l'accumulation pèsent aussi, mais elles sont
bornées. Elles mesurent surtout une incertitude, que la donnée résorbera, c'est-à-dire le registre DORA,
plus sûrement qu'une action directe (figure~\ref{fig:kpi}).

\begin{figure}[H]\centering
\figover{Z13_kpi_pilotables.png}
\caption{Des indicateurs descriptifs aux indicateurs pilotables. (a)~Écart de capital DORA attribué
par indicateur, chaque canal étant isolé (bleu : calibrable ; gris : borné), avec l'interaction qui
referme la cascade. (b)~Tableau de bord de remédiation, indicateurs classés par capital libéré.}
\label{fig:kpi}
\end{figure}

\subsection{Retour sur investissement de la conformité}
\label{sec:roi}
% SOURCES-SCRIPTS: 50 58 83

Le tableau de bord se lit en euros, et deux mots demandent une définition. Le portage est le coût de
détention du capital : sous Solvabilité~II, il se lit dans la marge de risque, un pourcentage annuel du
capital immobilisé. Au taux de $6\,\%$ du règlement délégué, les $14\,139$~M\euro{} que libère, pour le périmètre financier étudié, la
conformité totale valent $848$~M\euro{} par an, et non une fois. La fréquence en porte à elle seule
près de la moitié.

Le bénéfice a deux composantes, et ce calcul n'en monétise qu'une. Le portage économisé représente
$848$~M\euro/an, la sinistralité évitée $3\,247$, la charge annuelle moyenne passant de $3\,869$ à
$622$~M\euro. La composante laissée de côté vaut $3{,}8$ fois celle qui est retenue, de sorte que le
portage seul minore le bénéfice. Le retour de $6$ à $35$ ans est donc un majorant, pour un coût de
remédiation de $25$ à $150$~M\euro{} par grande entité, soit $5$ à $30$~Md\euro{} pour deux cents
entités.

Ajouter les pertes évitées raccourcit cette durée sous les mêmes hypothèses, sans borner le retour réel. Les deux composantes réunies donnent
$4\,095$~M\euro/an et un retour de $1$ à $7$ ans. Cette addition mêle un flux de compte de résultat à
un coût du capital. Cette convention de retour sur investissement n'est pas une sortie du modèle. Le chiffre de
référence reste le portage seul, avec le sens de sa borne.

Ce taux de convention n'est ni le coût moyen pondéré du capital de l'entité ni son taux
d'exigence sur un projet ; la durée calculée ne valide pas une décision d'investissement.

Le taux de $6\,\%$ est celui du règlement délégué~2015/35. La directive~(UE)~2025/2 le ramène à
$4{,}75\,\%$, taux qu'emploient les calculs d'échelle d'entité. Le portage monétisé tombe alors à
$672$~M\euro/an, soit une baisse de $21\,\%$, et le majorant du retour passe de $35$ à $45$ ans. Aucune
conclusion ne bascule : la durée reste un majorant, et la composante dominante du bénéfice n'est pas
monétisée. La chaîne publiée reste au taux historique, la calibration étant gelée
(figure~\ref{fig:roi}). La conformité n'est donc pas un simple coût. Elle libère du capital, que le modèle chiffre, face à un
coût externe et incertain, affiché comme tel. On priorise les leviers calibrables.

\begin{figure}[H]\centering
\figover{J6_roi_conformite.png}
\caption{Retour sur investissement de la conformité. (a)~Capital libéré par levier, classé, avec sa
valeur annuelle à $6\,\%$ (bleu : calibrable ; gris : borné). (b)~Coût de remédiation ponctuel contre
économie récurrente de portage du capital. Le retour sur le portage seul est de $6$ à $35$ ans, et
c'est un majorant, la perte évitée n'étant pas monétisée ici.}
\label{fig:roi}
\end{figure}

\section{Piliers laissés en perspective}
\label{sec:perspectives-pilier}

Deux piliers appelleraient un traitement propre, mais aucune donnée ne le soutient. Les modéliser par
des paramètres non identifiables contredirait la démarche, et nous les laissons en perspective.
\begin{itemize}\itemsep2pt
\item P1 (gouvernance) ressemble moins à un n\oe{}ud qu'à un modificateur systémique : une gouvernance
      défaillante charge la sensibilité $\rho$ de tous les piliers à la fois. Aucune donnée ne relie la
      qualité de gouvernance à une hausse du hasard, et cette relation est à borner plutôt qu'à
      estimer.
\item P5 (partage d'informations) agit plutôt comme un réducteur de fréquence, l'alerte précoce
      diminuant les entrées. Rien ne le mesure non plus.
\end{itemize}
Les diagnostics motivent des adaptations distinctes : hétérogénéité des pertes, choc commun et
détection. Ils n'identifient pas tous les paramètres correspondants. Les scénarios examinés
relèvent le niveau du capital sans renverser le classement des piliers. Leur usage reste
conditionnel aux hypothèses de chaque calcul ; un registre au niveau de l'entité serait nécessaire
pour estimer les mécanismes encore posés.
